\documentclass[10pt,a4paper]{amsart}
\usepackage[margin=19mm]{geometry}
\usepackage{amsmath,amssymb,mathtools}
\usepackage[scr=rsfs,cal=euler]{mathalfa}
\usepackage{xfrac}
\usepackage[unicode,hidelinks,bookmarks=false]{hyperref}
\newcommand{\ie}{i.e.\ }
\newcommand{\cf}{cf.\ }
\newcommand{\resp}{respectively\ }
\newcommand{\Subsection}{\subsection}
\providecommand{\nobreakdash}{\nobreak\hbox{-}}
\setlength{\emergencystretch}{2em}
\multlinegap0pt
\usepackage{amssymb}
\usepackage{algorithm}\usepackage{algpseudocode}
\usepackage{enumerate}
\usepackage{tikz}
\newtheorem{theorem}{Theorem}
\newcommand{\CONX}[1]{\mathcal{C}(#1)}
\newcommand{\transpose}[1]{{#1}^{\top}}
\title{Hardness of some optimization problems over correlation polyhedra}
\author{Alberto Caprara, Fabio Furini, Claudio Gentile, Leo Liberti, Andrea Lodi}
\date{Source: arXiv:2605.02896v1 (CC BY 4.0)}
\begin{document}
\maketitle

\noindent\emph{Source and licence.} Hardness of some optimization problems over correlation polyhedra. Version arXiv:2605.02896v1 (CC BY 4.0), distributed by arXiv under the Creative Commons Attribution 4.0 International licence, http://creativecommons.org/licenses/by/4.0/, as recorded in the arXiv licence metadata for this exact version and shown on the abstract page https://arxiv.org/abs/2605.02896v1. Full licence notice: \texttt{licenses/arxiv-2605.02896-CC-BY-4.0.txt}.\par\smallskip

\noindent\textit{Editorial scope.} Counted unit: the article's theorem thm:rkhard (source0.tex lines 468--471). The article's complete original proof of it is delivered with it (source0.tex lines 472--518, one \texttt{proof} environment of 47 lines). No mathematics is written, repaired or re-derived: the statement and the proof are the article's own lines, in the article's own notation, and no retained line is substituted or re-flowed.  Retained uncounted as the source-local context the proof needs: lines 455--464.  Nothing was omitted from any retained span, so the package carries no omission record.  No retained line contains a citation, so the wrapper carries no bibliography and no bibliography entry.  No retained line contains a figure environment, an image-inclusion command or drawing code, so no image is delivered and the package declares no asset dependency.  Editorial formatting only, in addition: an AMS \texttt{amsart} preamble that restates the article's own declarations and macros used by the retained lines, namely the environments \texttt{theorem} on separate counters, together with the standard packages those lines need.  The retained environments print in this document as Theorem 1 in order of appearance, whereas the article prints the same items with the numbering of its own full document.  The complete native source of the article as distributed in the arXiv e-print for this exact version is bundled unchanged as \texttt{sources/199756/source0.tex}.
\subsection{Rank in $\CONX{n}$}
\label{s:rank}
In this section, we prove the \textbf{NP}-hardness of the rank problem associated to $\CONX{n}$.
\begin{quote}
  Given a positive integer $q$ and a symmetric $n\times n$ matrix $\Gamma$, find a vector $p\in\mathbb{R}^{P_n}_+$ such that
  \begin{eqnarray}
    \|p\|_0 &\le& q \label{eq:rk1} \\
    \Gamma &=& \sum_{k\le P_n} p_k x_k\transpose{x}_k. \label{eq:rk2}
\end{eqnarray}
\end{quote}

\begin{theorem}
  The rank problem for $\CONX{n}$ is \textbf{NP}-hard.
  \label{thm:rkhard}
\end{theorem}

\begin{proof}
  Given an instance of \textsc{linear}-X3C with $|U|=3q$ for some integer $q\ge 1$, we assume without loss of generality that $U=\{1,\ldots,n-1\}$, so $n=3q+1$. We construct $\Gamma$ as follows:
  \begin{itemize}
  \item for each $i<n$: $\Gamma_{in}=\Gamma_{ni} = 1$;
  \item for each $i<j<n$:
    $\Gamma_{ij} = \Gamma_{ji} = \left\{\begin{array}{ll} 1 & \mbox{if $\exists S\in\mathcal{S}$ with $\{i,j\}\in S$,} \\
      0 & \mbox{otherwise;}
    \end{array}\right.
    $
  \item for $i\in U$: $\Gamma_{ii} = 1$;
  \item $\Gamma_{nn}=q$.
  \end{itemize}
  
  First, we assume that the \textsc{linear}-X3C instance is YES. So, there exists an exact cover $\mathcal{T}\subset\mathcal{S}$ that partitions $U$: for every $i<n$ there is a unique $S\in\mathcal{T}$ such that $i\in S$. Since $|U|=3q$ and it has a partition into triplets, the number of triplets in the partition is $q$. So we can write $\mathcal{T}=\{T_1,\ldots,T_q\}$, where each $T_k\in\mathcal{T}$ is equal to some $S_j\in\mathcal{S}$. For each $k\le q$ we define $x_k\in\{0,1\}^n$ as follows:
  \[
  x_{kn}=1 \qquad \mbox{and} \qquad \forall i<n\quad x_{ki}=\left\{\begin{array}{ll} 1 & \mbox{if $i\in T_k$,} \\ 0 & \mbox{otherwise.}\end{array}\right.
  \]
  Moreover, we let $p_k=1$ for all $k\le q$. Let $\bar{\Gamma}=\sum_{k\le q} p_k x_k\transpose{x}_k$: we claim that $\bar{\Gamma}=\Gamma$. For each $i<n$ exactly one triplet $T_k$ in $\mathcal{T}$ contains $i$, hence exactly one boolean vector $x_k$ has $x_{ki}=x_{kn}=1$, which implies $\bar{\Gamma}_{in}=1=\Gamma_{in}$. The same consideration also implies that $\bar{\Gamma}_{ii}=\Gamma_{ii}$ for all $i<n$. For $i<j<n$ we have $\bar{\Gamma}_{ij}=1$ iff $i,j$ both belong to a given triplet $T_k$ (for some $k\le q$). By definition of \textsc{linear}-X3C, $T_k$ must be the only triplet to contain both $i,j$: hence $\bar{\Gamma}_{ij}=\Gamma_{ij}$. Therefore, $\Gamma$ satisfies Eq.~\eqref{eq:rk2} with exactly $q$ nonzero coefficients, whence the rank of $\Gamma$ is $\le q$.

  Next, we assume that $(q,\Gamma)$ is a YES instance of the rank problem for $\CONX{n}$, i.e.~$\Gamma$  has rank $r\le q$, whence $\Gamma=\sum_{k\le r} p_k x_k\transpose{x}_k$ with $p>0$. For each $k\le r$ we define the support sets $R_k=\{i<n\;|\;x_{ki}=1\}$ of each $x_k$. For each $i<n$ we have
  \begin{equation}
    1=\Gamma_{ii}=\sum_{k\le r} p_k x_{ki}^2=\sum_{k\le r\atop i\in R_k} p_k.
    \label{eq:ii}
  \end{equation}
  Moreover, since $\Gamma_{ni}=\Gamma_{in}=1$ for all $i<n$, we have
  \[1=\Gamma_{in}=\sum_{k\le r} p_k x_{kn}x_{ki}=\sum_{k\le r\atop i,n\in R_k} p_k.\]
  Subtracting one equation from the other, we infer that any $x_k$ such that $x_{ki}=1$ must also have $x_{kn}=1$, otherwise the two sums could not both  be equal to $1$. Hence, for any $k\le r$, if $R_k$ contains $i<n$ it also contains $n$. Now, for any $k\le r$ consider $i<j<n\in R_k$: the $k$-th term contributes $p_k>0$ to $\Gamma_{ij}$. Since $\Gamma\in\{0,1\}^{n\times n}$ and $p\ge 0$, $\Gamma_{ij}=1$. By construction of $\Gamma$, there must be $S_h\in\mathcal{S}$ such that $i,j\in S_h$. By definition of \textsc{linear}-X3C, all pairs in $R_k$ must be in a single triplet, so $R_k\subseteq S_h$, implying $|R_k|\le 3$.

  At this point, we claim that $r=q$, all $p_k=1$, and $|R_k|=3$. We sum over the last column of $\Gamma$ and obtain $\sum_{i<n}\Gamma_{in}=3q$, since $\Gamma_{in}=1$ for all $i<n$, and $|U|=|\{i\;|\;i<n\}|=3q$. We also have
  \[\sum_{i<n}\Gamma_{in}=\sum_{i<n}\sum_{k\le r}p_k x_{ki} x_{kn}=\sum_{k\le r}p_k|R_k| x_{kn}=\sum_{k\le r}p_k|R_k|\]
  by definition of $\Gamma$, by definition of $R_k$, and because $x_{kn}=1$ whenever $R_k\not=\varnothing$. Note that $|R_k|\le 3$ implies
  \begin{equation}
    3q = \sum_{k\le r} p_k|R_k|\le 3\sum_{k\le r}p_k.
    \label{eq:3q}
  \end{equation}
  Now, by $\Gamma_{nn}=q$ we have
  \begin{equation}
    q = \Gamma_{nn}=\sum_{k\le r} p_k x_{kn}^2 =\sum_{k\le r} p_k x_{kn}=\sum_{k\le r} p_k,
    \label{eq:qq}
  \end{equation}
  again by definition of $\Gamma$, by $x_{kn}\in\{0,1\}$, and because $x_{kn}=1$ whenever $p_k>0$. Combining Eq.~\eqref{eq:3q}-\eqref{eq:qq} we obtain
  \[3q\le 3\sum_{k\le r} p_k = 3q,\]
  whence equality must hold in Eq.~\eqref{eq:3q}. Therefore $|R_k|=3$ for all $k$ with $p_k>0$. To settle the claim, we remark the following facts: (i) $|R_k|=3$ for all $k\le r$, (ii) $\sum_{k\le r} p_k=q$, (iii) $r\le q$, and (iv) $p_k>0$ for all $k\le r$. These imply that $r=q$ and $p_k=1$ for all $k\le r$, as claimed.
  %we sum Eq.~\eqref{eq:ii} over all $i<n$, which yields: \[ 3q=\sum_{i<n} 1 = \sum_{i<n}\sum_{k\le r\atop i\in R_k}p_k=\sum_{k\le r} p_k|R_k|=\sum_{k\le r} 3p_k=3\sum_{k\le r} p_k=3q\]
  
Finally, by Eq.~\eqref{eq:ii} and the fact that $p_k=1$ for all $k\le r$, for each $i<n$ we have \[1=\sum_{k\le r\atop i\in R_k}1,\] so each $i<n$ belongs to exactly one $R_k$. Therefore, $R_1,\ldots R_r$ form a partition of $U$ into $q$ disjoint triplets: then $\mathcal{T}=\{R_1,\ldots R_r\}$ is an exact cover of $U$, which implies that the \textsc{linear}-X3C is a YES instance, as claimed.
\end{proof}

\end{document}
